\documentclass[10pt,a4paper]{extarticle}
\usepackage[utf8]{inputenc}
\usepackage[T1]{fontenc}
\usepackage[spanish]{babel}
\usepackage[a4paper,margin=1.25cm,headheight=12pt]{geometry}
\usepackage{tabularx,array,booktabs}
\usepackage{enumitem}
\usepackage[hidelinks]{hyperref}
\usepackage{fancyhdr}
\setlength{\parindent}{0pt}
\setlength{\parskip}{2pt}
\setlist{nosep,leftmargin=1.3em}
\renewcommand{\arraystretch}{1.12}
\pagestyle{fancy}
\fancyhf{}
\fancyhead[L]{CGB Academy · Ingeniería de Software I}
\fancyhead[R]{Actividad de necesidades y requisitos}
\fancyfoot[C]{\thepage}
\newcolumntype{Y}{>{\raggedright\arraybackslash}X}
\newcolumntype{P}[1]{>{\raggedright\arraybackslash}p{#1}}
\begin{document}
\small
\begin{center}
{\Large\bfseries Identificación y validación de necesidades}\par
Proyecto CGB Academy · Informe breve · 21 de octubre de 2026
\end{center}
\vspace{-4pt}
\textbf{Estado de evidencia:} el equipo reporta conversatorios reales con administración de ventas, TI y operaciones académicas. Las fotos recibidas muestran sesiones de trabajo, pero no contienen respuestas. Las necesidades aquí se trazan a documentos; los testimonios aún no están transcritos por rol.

Base consultada: especificación funcional y no funcional de CGB Academy, Jira CGB-51/CGB-56, \texttt{src/app/colaborador/page.tsx} y manual local de seguimiento comercial por WhatsApp. El registro detallado está en \texttt{../registro\_conversatorios.md}.

\section*{A. Stakeholders y persona}
\footnotesize
\begin{tabularx}{\linewidth}{P{2.8cm} Y P{1.5cm} P{1.5cm} P{1.5cm}}
\toprule
Stakeholder & Expectativa documentada o por validar & Poder & Interés & Prioridad \\
\midrule
Administración de ventas & Dar seguimiento a prospectos por estado y canal & Medio & Alto & Alta para consulta \\
\addlinespace[2pt]
Área de TI & Mantener acceso, credenciales, datos e infraestructura seguros & Alto & Alto & Alta \\
\addlinespace[2pt]
Operaciones académicas & Gestionar capacitaciones, rutas y seguimiento & Alto & Alto & Alta \\
\addlinespace[2pt]
Colaborador interno & Completar inducción y saber qué sigue & Medio & Alto & Alta \\
\addlinespace[2pt]
Público (estudiantes/docentes) & Encontrar orientación pública sin iniciar sesión & Bajo & Alto & Media \\
\bottomrule
\end{tabularx}
\normalsize
Se consulta de cerca a TI y operaciones académicas por su poder e interés altos, y a ventas y al colaborador por su interés alto y participación directa. El público se mantiene informado por su interés alto y poder bajo. Aunque ventas es stakeholder prioritario para consulta, la venta de cursos sigue excluida; cualquier cambio requiere decisión del Product Owner. Los niveles son estimaciones documentales. \textbf{Persona anonimizada:} participante real de operaciones académicas, según el equipo. Su objetivo documentado es gestionar capacitaciones, rutas y seguimiento. Frustraciones, dispositivo y cita textual no constan en las fuentes; se omiten datos personales.

La matriz poder--interés es provisional: alto/alto se gestiona de cerca; bajo/alto se mantiene informado. Las fotos aportan contexto de las sesiones, no respuestas ni vínculo entre rostro y rol.

\textbf{Elicitación y técnica adicional:} el equipo informa que conversó con los tres roles. Se analizaron la especificación y el manual comercial; esta técnica identifica necesidades candidatas, no reemplaza testimonios. Preguntas de seguimiento para operaciones académicas: (1) ¿Cómo se crea y asigna una capacitación reciente? (2) ¿Cómo se determina la ruta según área o puesto? (3) ¿Qué pasos se repiten o se resuelven fuera del sistema? (4) ¿Cómo se comprueba que se terminó la versión correcta? (5) ¿Quién debe ver el avance y qué datos deben ser privados? No se afirma que estas preguntas se formularon en las sesiones.

\section*{B. Requisitos revisados}
\footnotesize
\begin{tabularx}{\linewidth}{P{2.8cm} P{3.5cm} Y}
\toprule
Tipo e inicial & Problema & Reformulación y verificación \\
\midrule
\textbf{Funcional — RF-011:} panel con capacitaciones y porcentaje ``en tiempo real''. & ``Tiempo real'' y estados sin definir. & Mi Onboarding muestra ruta ordenada y estados Pendiente/En progreso/Completada. Para N capacitaciones y M completadas calcula $\mathrm{round}(100\times M/N)$; si N=0 muestra 0\%. \textbf{Verificar:} pruebas de ruta vacía y estados mixtos; comparar lista y porcentaje. \\
\textbf{No funcional — RNF-004:} carga inicial $<1{,}5$ s en 4G. & No fija pantalla, perfil de red ni medición. & En build de producción, el contenido principal de Mi Onboarding tendrá LCP p95 $<1{,}5$ s en 20 mediciones con dispositivo, perfil 4G y datos acordados. \textbf{Verificar:} registrar parámetros y las 20 medidas; el PO confirma el umbral. \\
\textbf{Regla de negocio — CGB-56:} capacitación pasa a Completada al finalizar. & No define si basta llegar al último slide ni cómo se desbloquea el siguiente módulo. & Se marca Completada solo tras recorrer la secuencia y confirmar Finalizar en el último slide; el módulo siguiente se habilita si sus prerrequisitos están completados. \textbf{Verificar:} BDD para avance parcial, confirmación final y desbloqueo. Pendiente de aprobar con el PO. \\
\textbf{Restricción MVP:} certificados excluidos. & CGB-51 incluía constancia descargable, fuera de alcance. & El MVP no genera insignias ni certificados; al completar solo muestra el estado de ruta. \textbf{Verificar:} revisar criterios Jira y que la interfaz no ofrezca descarga. \\
\bottomrule
\end{tabularx}
\normalsize

\newpage
\section*{C. Backlog refinado y trazabilidad}
\textbf{Necesidad N2 — retomar y confiar en el guardado.} \textbf{\href{https://cgb-academy.atlassian.net/browse/CGB-56}{CGB-56 · HU-012}} (prioridad \textbf{Highest}; ``Crítica'' en la especificación). Como colaborador, quiero guardar y retomar el último slide confirmado para continuar sin perder avance y reducir repetición de orientación (impacto esperado, pendiente de validar). Stakeholder: colaborador interno, interés alto y poder medio estimados. Trazada a RF-012; falta vincularla a una respuesta específica.
\begin{itemize}
\item \textbf{Aceptación:} guardar el slide confirmado; reanudar allí; completar tras confirmar todos los slides y pulsar Finalizar en el último; desbloquear la siguiente capacitación al cumplir prerrequisitos; bloquear lectura/escritura entre cuentas; repetir una confirmación no duplica avance.
\item \textbf{Tareas Jira existentes:} CGB-57 pruebas; CGB-58 endpoint y persistencia; CGB-59 idempotencia y UX; CGB-60 BDD y PR a \texttt{main}.
\item \textbf{AG-C01:} progreso privado y aviso claro de guardado; pruebas de autorización con dos cuentas. El beneficio de reducir repetición queda por validar.
\end{itemize}

\textbf{Necesidad N1 — conocer lo asignado y el siguiente paso.} \textbf{\href{https://cgb-academy.atlassian.net/browse/CGB-51}{CGB-51 · HU-011}} (prioridad \textbf{Highest}; ``Crítica'' en la especificación), consume el avance de CGB-56. Como colaborador autenticado, quiero ver mi ruta, estados y porcentaje para saber qué me falta y reducir consultas manuales (impacto esperado, pendiente de validar). Evidencia documental: RF-011 y Jira; falta vincularla a una respuesta específica.
\begin{itemize}
\item \textbf{Aceptación:} lista ordenada y estados visibles; porcentaje según fórmula de pág. 1; continuar desde el punto guardado; aislar datos por cuenta y explicar el guardado; no ofrecer certificado fuera del MVP.
\item \textbf{Tareas Jira existentes:} CGB-52 pruebas; CGB-53 API autenticada; CGB-54 panel accesible; CGB-55 BDD y PR a \texttt{main}.
\item \textbf{AG-C01:} no autenticados no reciben datos y cada cuenta consulta solo su avance; verificar con dos cuentas.
\end{itemize}

\textbf{Nota Jira (revisado el 21/10):} CGB-51 y CGB-56 siguen en Sprint 2, estado ``Tareas por hacer'', prioridad Highest y etiqueta \texttt{evidence-pending}. Se actualizó solo la sección de evidencia para registrar los conversatorios reportados y aclarar que sus respuestas aún no validan N1/N2. No se cambiaron estado, sprint, prioridad ni responsables. CGB no ofrece tipo ``Historia''; se conservan como ``Tarea'' con formato de historia de usuario.

\section*{Reflexión y pendientes}
Pasamos por alto que el colaborador necesita confiar en el guardado y privacidad del avance, y que ventas ya sigue prospectos por WhatsApp. Ventas queda fuera del MVP hasta decisión del Product Owner. Antes de cambiar CGB-51/CGB-56, hay que registrar por rol las respuestas de los conversatorios y contrastar criterios con operaciones académicas y TI. También falta confirmar el umbral de rendimiento y aprobar la regla de visibilidad.
\end{document}
